\documentclass[10pt]{article}
\usepackage[margin=0.9in]{geometry}
\usepackage{amsmath,amssymb,amsthm,booktabs,hyperref,microtype}
\setlength{\emergencystretch}{1em}
\newtheorem{theorem}{Theorem}
\newtheorem{lemma}{Lemma}
\newcommand{\E}{\mathbb E}
\newcommand{\Prb}{\mathbb P}
\newcommand{\Var}{\operatorname{Var}}
\title{A Single-Run Refinement of the Support-Testing Constants}
\author{}\date{October 22, 2026}
\begin{document}\maketitle
\begin{abstract}
Keeping the statistic and all polynomial parameters unchanged, we strengthen
the concentration step of the accompanying exact-certificate analysis.
The integer pair $(C_0,L_0)=(70,14)$ has uniform single-run failure below
$0.0009$, and $(70,16)$ has failure below $0.000057$.
Thus the previously sufficient 27 repetitions of $(70,16)$ are unnecessary
for a one-percent guarantee. The real pair $(144(\log2)^2,14)$ is also certified.
An asymptotic null construction also forces $C\ge69.00627138\ldots$ for
this unchanged tester, proving that 70 is its smallest feasible integer
sampling constant for uniform one-percent error. The best real constant
and smallest branch threshold are not determined.
The proof includes the infinite range of input sizes; simulation does not
establish the theorem. It builds on the unchanged expectation, coefficient,
and interval lemmas in the repository's previous full manuscript.
\end{abstract}

\section{Construction and exact statement}
Let $p$ be any finite or countably supported discrete distribution. For an
integer $n\ge1$ and $0<\epsilon<1$, put
$L=1+\log n$, $m=Cn/(\epsilon L)$, and draw independent
$N_i\sim\operatorname{Poi}(mp_i)$. Write $D=\sum_i1_{N_i>0}$,
$S=|\operatorname{supp}p|$, $\Delta_r=\sum_{i>r}p_{(i)}$, and
$e_\epsilon=\min\{r\ge1:\Delta_r\le\epsilon\}$.
For $L<L_0$ use $\widetilde S=D$. Otherwise use exactly the previous construction:
\[
 a_0=\epsilon/(8n),\quad b=64L/m,\quad
 J=\lfloor\log_2(256L^2/C)\rfloor,\quad a_j=2^ja_0,\quad
 \delta_j=2^{-j-2},\quad k_j=\lceil\sqrt{b/a_j}(j+3)\log2\rceil.
\]
Set $P_j(x)=-\rho_jT_{k_j}((b+a_j-2x)/(b-a_j))=\sum_\ell c_{j,\ell}x^\ell$,
with $\rho_j=T_{k_j}((b+a_j)/(b-a_j))^{-1}$; put $f_j(0)=0$,
$f_j(\ell)=1+c_{j,\ell}\ell!/m^\ell$ for $1\le\ell\le k_j$, and
$f_j(\ell)=1$ otherwise. Then
\[
 Z_j=\sum_i f_j(N_i),\qquad
 \widetilde S=\max\{D,\max_j(Z_j-\delta_j n)/(1+\delta_j)\}.
\]
The binary decision is Yes iff $\widetilde S\le n$; the integer output is
$\widehat S=\min\{n,\max\{1,\lceil\widetilde S\rceil\}\}$.
There is no budget ceiling or degree truncation in this theorem.

\begin{theorem}
For every row below, uniformly over $p,n,\epsilon$, the promised binary test
has error less than $B$, and
$\Prb\{\min(e_\epsilon,n)\le\widehat S\le S\}>1-B$.
On either binary promise these two outputs are simultaneously correct with
probability greater than $1-B$.
The promises are $S\le n$ and $\Delta_n\ge\epsilon$; no binary correctness
label is assigned in the intervening gap.
\begin{center}
\begin{tabular}{rrr}
\toprule $C$&$L_0$&exact rational upper bound $B$\\\midrule
$144(\log2)^2$&14&0.0015\\
$144(\log2)^2$&15&0.00038\\
$144(\log2)^2$&16&0.00031\\
$144(\log2)^2$&18&0.00028\\
70&14&0.0009\\70&15&0.000081\\70&16&0.000057\\70&18&0.000040\\
\bottomrule
\end{tabular}
\end{center}
\end{theorem}

\section{What is inherited, and what changes}
The unchanged analytic ingredients are proved in the
\href{../artifacts/joint_constants_submission/submission_proof.pdf}{earlier exact-certificate manuscript},
especially its expectation, integral/range, individual-coefficient, and
variance lemmas. They give $\E f_j(N_i)\le1+\rho_j$,
$\rho_j\le\delta_j^2/2$, an envelope $F_j\ge\|f_j\|_\infty$, and
a factorial-coefficient sum $I_j$ which bounds the range length of $f_j$.
The certificate uses $V_j=I_j^2/4$, improved at scale zero to
$V_0=\min(I_0^2/4,F_0^2)$ using the individual-coefficient envelope.
Thus variance is at most $V_j$, and squared range is at most $4V_j$. Also
\[
 \sum_i\Var f_j(N_i)\le F_j^2\E D.
\]
The latter follows from $f_j(0)=0$ and
$\Var f_j(N_i)\le\E f_j(N_i)^2\le F_j^2\Prb(N_i>0)$.
Only concentration inequalities and the cutoff in the occupancy case split
are strengthened here. Neither the sampling rule nor the estimator changes.

The original parameter lemma extends to $L\ge14$: for $69<C\le70$,
$256L^2/C>1$, $64L^2/(Ce^{L-1})<1$, and
\[
 k_j/L\le\sqrt{512/C}\,3\log2+1/14<6<64/\sqrt2.
\]
The first two expressions are monotone on this range. The inequalities at
14 are certified using outward rational intervals, so all expectation
lemmas continue to apply.

\section{Concentration lemmas with explicit constants}
For independent variables of range lengths $r_i$, Hoeffding's argument gives
\[
 \Prb\{\textstyle\sum_i(X_i-\E X_i)\ge x\}
 \le\exp\{-2x^2/\textstyle\sum_i r_i^2\}.
\]
For completeness, if $X\in[a,b]$, the second derivative of
$\log\E e^{tX}$ is a variance under exponential tilting, at most
$(b-a)^2/4$. Integrating twice at zero bounds the centered log moment
generating function by $t^2(b-a)^2/8$; exponential Markov and optimization
prove the display. The lower tail follows by negation.

We will also use the following Bernstein form: if independent centered
$Y_i\le B$ have total variance at most $\sigma^2$, then
\[
 \Prb\{\textstyle\sum_iY_i\ge x\}
 \le\exp\{-x^2/(2\sigma^2+2Bx/3)\}.
\]
One derivation uses
\[
 e^{ty}\le1+ty+\frac{y^2}{B^2}(e^{tB}-1-tB)\quad(y\le B,t\ge0),
\]
which follows by the integral remainder formula and monotonicity in $y$.
Since $k!\ge2\cdot3^{k-2}$ for $k\ge2$, the log moment generating
function is at most $\sigma^2t^2/[2(1-tB/3)]$. Choosing
$t=x/(\sigma^2+Bx/3)$ proves the stated bound.

An endpoint variant is useful for the upper tail. If $Y\in[-A,B]$,
$A,B>0$, and $\E Y\le0$, convexity bounds its moment generating function
by that of the mean-zero distribution on $\{-A,B\}$. This distribution
has variance $AB$ and upper endpoint $B$. Thus the same Bernstein bound
holds with variance proxy $AB$, even when $Y$ itself is not centered.

\section{Uniform upper and lower errors}
Define, as in the previous proof,
\[
 v_j=V_j/(n\delta_j^2),\quad w_j=F_j^2/(n\delta_j^2),\quad
 d_j=2-\delta_j/2,\quad\kappa_j=14-\delta_j/2-8\delta_j^2.
\]
For $S\le n$, the event $\widetilde S>S$ requires a scale's threshold
$(1+\delta_j)S+\delta_j n$ to be exceeded. Its gap above
$(1+\delta_j^2/2)S$ is at least
$x=\delta_j n+(\delta_j-\delta_j^2/2)S$.
For $0<S/n\le1$, the previous elementary maximization gives
$S/x^2\le1/(n\delta_j^2d_j^2)$, and $1/x\le1/(\delta_j n)$.

Writing $G_j=I_j-1$, positivity of absolute coefficients implies
$\min(0,1-G_j)\le f_j\le1+G_j$. Centering at
$1+\delta_j^2/2$ therefore permits the endpoint envelope
\[
 A_j=1+\delta_j^2/2+\max(0,G_j-1),\qquad B_j=G_j.
\]
Let
\[
 U=\sum_j\min\left\{
 \frac{v_j}{v_j+d_j^2},\ e^{-d_j^2/(2v_j)},\
 \exp\left[-\left(\frac{2A_jG_j}{n\delta_j^2d_j^2}
                  +\frac{2G_j}{3\delta_j n}\right)^{-1}\right]\right\}.
\]
The three terms are respectively the old Cantelli bound, the range-based
Hoeffding bound, and the endpoint Bernstein bound just proved. Their minimum
is valid because each bounds the same event. A union bound gives
$\Prb(\widetilde S>S)\le U$.

For the lower error fix $r\le n$ with $\Delta_r\ge\epsilon$.
Use the previous three-case expectation argument, whose heavy-category
case has failure at most $H(L)=e^{-31L-1}+e^{-C}$ and whose other two cases
have rejection margin at least $\kappa_j\delta_j n$ at some scale $j$.
Change the occupancy split to $\alpha=33/32$. If $\mu=\E D>\alpha n$,
independence of occupancy indicators and exponential Markov give
\[
 \Prb(D\le n)\le e^{-n(\alpha-1-\log\alpha)}
 \le e^{-n/2112}.
\]
Here $g(x)=x-1-\log x\ge(x-1)^2/(2x)$ for $x\ge1$; subtracting the
right side from $g$ gives derivative $(1-1/x)^2/2\ge0$.
If $\mu\le\alpha n$, the variance is at most $\alpha n F_j^2$ and
the centered summands in the lower tail have upper bound $2F_j$.
Consequently a valid general lower-error bound is
\[
 R=\max\left\{H(L),e^{-n/2112},
 \max_j\min\left[\frac{\alpha w_j}{\alpha w_j+\kappa_j^2},
 \exp\left(-\frac{\kappa_j^2}{2\alpha w_j+(4/3)\kappa_j\sqrt{w_j/n}}\right)
 \right]\right\}.
\]
For $S\le n$, the finite-support variance and Hoeffding bounds instead give
\[
 R_f=\max\left\{H(L),\max_j\min\left[
 \frac{v_j}{v_j+\kappa_j^2},e^{-\kappa_j^2/(2v_j)}\right]\right\}.
\]
All infinite sums above are legitimate: $\E D\le m<\infty$,
$\sum_i\E|f_j(N_i)|\le F_j\E D$, and the centered series converges in
$L^2$. The finite-sum exponential bounds pass to the limit by Fatou's
lemma. No finite-alphabet assumption is added.

The common bound is $B(L)=\max\{U+R_f,R\}$. For $S\le n$, the event
$\widetilde S\le S$ and $\widetilde S>\min(e_\epsilon,n)-1$ (when the
minimum exceeds one) ensures both outputs are correct, so $U+R_f$ suffices.
For $S>n$ under the alternative, rejection $\widetilde S>n$ ensures both
outputs are correct; $R$ suffices. In the gap only the support bound is
required and the same $R$ argument applies. This proves the simultaneous
statement, without adding two separate failure bounds.

\section{Exact verification on all input sizes}
On $14\le L\le40$, the interval envelopes from the previous manuscript are
used on all 260 closed slabs $[i/10,(i+1)/10]$. Each slab includes every
possibly active scale and upper bounds every degree ceiling; there is no
unchecked gap at a scale or degree transition. The new probability formulas
are increasing in their variance, coefficient, and inverse-$n$ arguments,
so substituting their slab upper bounds is sound.

On $L\ge H=40$ the previous half-line proof is unchanged except for the
choice of $H$. For scales $j\le7$, its decreasing envelopes bound $w_j$;
we conservatively use $v_j\le w_j$ and omit the endpoint Bernstein option
for the upper tail. At scale zero the condition is exactly
$144(\log2)^2/C\le1$, including symbolic equality for $C=C_*$.
The other derivative signs and $H\ge14$ are checked exactly.
For $j\ge8$ the old decreasing sum envelope $T(H)$ bounds
$\sum_j w_j$. These omitted scales contribute at most $T(H)$ to $U$,
$T(H)/81$ to $R_f$, and $\alpha T(H)/81$ to $R$, using Cantelli and
$d_j>1$, $\kappa_j>9$. The maximum inverse-$n$ is $e^{1-H}$.
Thus the entire unbounded half-line is certified, not sampled on a grid.

In the counting branch the discovery coupling from the previous proof gives
$\Prb(D\le r)\le\Prb\{\operatorname{Poi}(Cn/L)\le n\}$.
Partition $1\le L\le L_0$ into intervals of width $1/4$. On a slab $[a,b]$,
use $n\ge e^{a-1}$ and $x=C/L\ge C/b>1$ to bound this probability by
\[
 \exp\{-e^{a-1}(C/b-1)^2/(2C/b)\}.
\]
Every such bound is certified below $10^{-10}$ for the listed pairs.
The counting branch never overestimates $S$.

The certificate uses only the previous 160-bit dyadic interval library,
integer operations, and rational operations; its logarithm, square-root,
exponential, and $\pi$ enclosures are proved in the previous arithmetic
appendix. For extremely negative exponents we safely enclose
$e^x$ in $[0,2^{-128}]$ when $x\le-128$, using $e>2$.
This avoids enormous intermediate integers and does not round a probability
down to zero. All asserted rational targets are strictly larger than the
certified upper endpoints. A separate 512-bit Arb implementation recomputes
the envelopes and probability expressions and audits the saved data.

\section{A necessary constant: ruling out 69}
This section concerns the unchanged tester above, not all possible support
testers. Both $C$ and $L_0>1$ are fixed finite constants.
Define
\[
 a(C)=\sqrt{512/C}\,3\log2,\qquad
 \gamma(C)=a(C)\operatorname{arsinh}(a(C)/32)
                -\sqrt{1024+a(C)^2}+32.
\]
The function $\gamma$ is continuous and strictly decreasing in $C>0$:
its derivative with respect to $a$ is $\operatorname{arsinh}(a/32)>0$,
and $a$ decreases in $C$. It decreases from infinity to zero, so there is
a unique $C_{\rm crit}$ such that $\gamma(C_{\rm crit})=1/2$.

\begin{theorem}[Necessary constant for the unchanged construction]
If $1\le C<C_{\rm crit}$, there are null distributions $p_n$ with support
exactly $n$ and $\epsilon_n\in(0,1)$ such that the false-rejection probability
has limit inferior at least $1/2$ as $n\to\infty$.
If $0<C\le1$, a counting-branch alternative has error above $1/2$.
These obstructions persist under any fixed odd number of independent runs
with majority aggregation. Moreover,
\[
 69.00627138<C_{\rm crit}<69.00627139.
\]
Consequently 70 is the smallest integer $C$ admitting a uniform one-percent
guarantee with any fixed finite $L_0>1$ and this unchanged single-run rule.
Among integer $C$ and fixed odd majority repetition counts, the minimal
feasible total sampling coefficient is also 70.
\end{theorem}

\begin{proof}
Fix $1\le C<C_{\rm crit}$. Only scale zero is needed. Its dimensionless
parameters are $A=C/(8L)$, $B=64L$, $k=\lceil a(C)L\rceil$.
Let $u_\ell=|c_{0,\ell}|\ell!/m^\ell$ and
$I_L=\sum_{\ell=0}^k u_\ell$, where $u_0=1$.
These quantities do not depend on $\epsilon$, because $A$ and $B$ do not.
The exact factorial-weight integral identity from the previous proof is
\[
 I_L=\int_0^\infty e^{-s}
 \frac{T_k(z_L+2s/(B-A))}{T_k(z_L)}\,ds,
 \qquad z_L=(B+A)/(B-A).
\]
The denominator tends to $\cosh(6\log2)$, since
$k\operatorname{arcosh}z_L\to6\log2$.
Writing $s=yL$, the logarithm of the integrand divided by $L$ tends,
uniformly on compact neighborhoods of every $y>0$, to
\[
 \phi(y)=-y+a(C)\operatorname{arcosh}(1+y/32).
\]
This follows directly from $T_k(z)=\cosh(k\operatorname{arcosh}z)$ for
$z>1$ and the limits of $k/L$, $(B-A)/L$, and $z_L$.
Differentiating $\phi$ shows that its positive maximizer is
$y_*=\sqrt{1024+a(C)^2}-32$, with maximum $\gamma(C)$.
For any $\eta>0$, integration over a sufficiently short fixed neighborhood
of $y_*$ therefore yields, for all sufficiently large $L$,
\[
 I_L\ge\exp\{(\gamma(C)-\eta)L\}.
\]
Choose $\eta$ so that $\gamma(C)-\eta>1/2$.
No unproved upper version of Laplace's method is needed here.

Choose $\ell_L\in\{1,\ldots,k\}$ maximizing $u_\ell$, so
$u_{\ell_L}\ge(I_L-1)/k$. Define the exact null distribution
\[
 p_n=(1/2,\underbrace{1/[2(n-1)],\ldots,1/[2(n-1)]}_{n-1\text{ entries}}),
 \qquad
 \epsilon_n=\frac{Cn}{2(n-1)\ell_L L}.
\]
Since $\ell_L\ge1$, the chosen $\epsilon_n$ is in $(0,1)$ eventually.
Its prescribed mean is exactly $m=2(n-1)\ell_L$, so each rare label has
Poisson mean $\ell_L$, and the head label has mean $(n-1)\ell_L$.
This choice is not circular: the weights determining $\ell_L$ depend only
on $C,L$, as noted above.

Let $X_L=f_0(\operatorname{Poi}(\ell_L))$, with mean $\mu_L$ and variance
$\sigma_L^2$. The expectation lemma applies for large $L$: for $C\ge1$,
$a(C)<64$, while $\sqrt{B(B-A)}/L\to64$; also the endpoint $b<1$
eventually. It implies $0\le\mu_L\le1+\rho_0\le33/32$.
Also $|X_L|\le I_L$.

For an integer $\ell\ge1$, the Poisson mass at $\ell$ is a mode. Chebyshev's
inequality gives probability at least $3/4$ to the interval of radius
$2\sqrt\ell$, containing at most $4\sqrt\ell+1$ integers. Thus
$\Prb\{\operatorname{Poi}(\ell)=\ell\}\ge3/(20\sqrt\ell)$.
At $\ell_L$, the alternating-sign formula for $f_0$ gives
$|f_0(\ell_L)|\ge u_{\ell_L}-1\ge I_L/(4k)$ eventually. Hence
\[
 \sigma_L^2\ge\frac{3I_L^2}{320k^{5/2}}-(33/32)^2
              \ge c\,\frac{I_L^2}{k^{5/2}}
\]
for an absolute positive $c$ and all sufficiently large $L$.

The $n-1$ rare-label scores are independent copies of $X_L$. Put
$s_n=\sqrt{n-1}\,\sigma_L$. Their largest possible centered magnitude
divided by $s_n$ is at most
\[
 \frac{2I_L}{\sqrt{n-1}\,\sigma_L}
       =O\left(\frac{k^{5/4}}{\sqrt n}\right)\longrightarrow0.
\]
Their normalized variances sum to one. Thus the Lindeberg condition is
eventually zero for each fixed cutoff, and the triangular-array central
limit theorem gives
\[
 \frac{\sum_{i=2}^n f_0(N_i)-(n-1)\mu_L}{s_n}
       \ \Longrightarrow\ \mathcal N(0,1).
\]
For a reference statement see Theorem 9 in
\href{https://people.eecs.berkeley.edu/~jordan/courses/210B-spring07/lectures/stat210b_lecture_2.pdf}
{Jordan--Kleiner, Berkeley Stat210B, Lecture 2 (2007)}.
The bounded normalized summands above explicitly verify its hypothesis.
The exponential lower bound on $I_L$ and $k=O(L)$ imply $n/s_n\to0$.
The head score equals 1 with probability tending to one, since its mean
$(n-1)\ell_L\ge n-1$ greatly exceeds $k=O(\log n)$.
It follows that
\[
 \Prb\{Z_0>3n/2\}\longrightarrow1/2.
\]
But $Z_0>3n/2$ alone implies $\widetilde S>n$, since $\delta_0=1/4$.
The polynomial branch is eventually active for every fixed finite $L_0$,
proving the claimed false-rejection obstruction.

For $0<C\le1$, take $n=1$, a uniform distribution on $K$ labels, and
$\epsilon=1-1/K$. This is an alternative and uses the counting branch
because $L_0>1$. As $K\to\infty$, collisions vanish and the error tends to
$\Prb\{\operatorname{Poi}(C)\le1\}=(1+C)e^{-C}\ge2/e>1/2$.
For fixed odd $r$, the binomial majority-error polynomial is increasing
and equals $1/2$ at single-run error $1/2$. Neither obstruction can be
removed by such repetition.

Finally, the root bracket is certified by
\texttt{parameter\_search/lower\_bound.py}. It writes
$\gamma(C)=32[x\operatorname{arsinh}x-\sqrt{1+x^2}+1]$ with
$x^2=(9/2)(\log2)^2/C$, and encloses $\operatorname{arsinh}x$ using the
alternating binomial series on $0<x<1$. Even and odd partial sums give
rational upper and lower bounds, respectively. The resulting sign tests
at both rational endpoints and at 69 are also independently recomputed
with 512-bit Arb. The positive theorem at $(70,14)$ completes the integer
optimality conclusion.
\end{proof}

\section{Scope of the optimization claim}
The exploratory search considers $C=C_*$ and every integer $70\le C\le150$,
with integer $10\le L_0\le30$, followed by exact verification of the listed
candidates. In this search and certificate family, $C_*$ has the lowest
sampling coefficient; 70 is the smallest covered integer. Fourteen is the
smallest integer threshold passing the explored bound at either of those
sampling constants. Failure of an upper-bound calculation at a smaller
parameter does not prove that the algorithm fails there.
In particular, $C\ge C_*$ is a limitation of the sufficient half-line
envelope. The separate necessary result above is $C\ge C_{\rm crit}$.
The best real sampling constant for this construction lies between
$69.00627138\ldots$ and $69.18523400\ldots$; the multiplicative gap is
less than $0.26\%$. Neither result is an information-theoretic bound over
all algorithms. The smallest feasible $L_0$ at $C=70$ remains unresolved;
14 is sufficient, not proved necessary. Finite empirical success for 69
does not contradict its asymptotic obstruction.

For implementation a useful compromise is $(70,14)$: it preserves a sizable
margin below one-percent error while activating the polynomial at
$n=442414$. Keeping $(70,16)$ activates it at $n=3269018$, has a smaller
certified error bound, and avoids polynomial computation in between.
Both use the same single-run expected sample count $70n/(\epsilon L)$.
Relative to the previously certified coefficient 1025 this is a reduction
of approximately $93.17\%$ in expected samples. It is not a claim of
the same factor in CPU time or simulation wall time. All statements assume
the iid Poisson model and exact statistic, with numerical certification
handled separately by the implementation.
\end{document}
